Un elemento radiactivo $A$ se transforma en $B$, que a su vez se transforma en $C$. Los ritmos de desintegración de primer orden son $k_1,k_2>0$, medidos en tiempo$^{-1}$. Inicialmente sólo hay $A$, en cantidad $x_0>0$. Sean $x(t)$ e $y(t)$ las cantidades de $A$ y $B$ para $t\ge0$. Use una unidad de cantidad compatible con la transformación uno a uno.

\begin{enumerate}[label=(\alph*),ref=\alph*,leftmargin=*]
\item Halle $y(t)$.
\item Determine la cantidad máxima de $B$ y el instante en que se alcanza.
\item La semivida de $A$ es el tiempo necesario para reducir su cantidad a la mitad. Demuestre que coincide con $t_*$ para una relación $k_1=ak_2$, y halle $a$.
\end{enumerate}
